\documentclass[11pt]{ltjsarticle}
\usepackage{luatexja-fontspec}
\usepackage{amsmath,amssymb}
\usepackage[margin=22mm]{geometry}
\usepackage{graphicx}
\usepackage{longtable,booktabs,array}
\usepackage{calc}
\usepackage{xcolor}
\usepackage{hyperref}
\hypersetup{colorlinks=true,linkcolor=blue,urlcolor=blue,citecolor=blue}
\makeatletter
\def\maxwidth{\ifdim\Gin@nat@width>\linewidth\linewidth\else\Gin@nat@width\fi}
\makeatother
\setkeys{Gin}{width=\maxwidth,keepaspectratio}
\setcounter{secnumdepth}{0}
\setlength{\parindent}{0pt}
\setlength{\parskip}{0.5em}
\providecommand{\tightlist}{\setlength{\itemsep}{0pt}\setlength{\parskip}{0pt}}
\providecommand{\pandocbounded}[1]{#1}
\providecommand{\real}[1]{#1}
\newcounter{none}
\begin{document}
\section{補遺：揺らぎの和ゼロから不確定性へ}\label{ux88dcux907aux63faux3089ux304eux306eux548cux30bcux30edux304bux3089ux4e0dux78baux5b9aux6027ux3078}

\subsection{------半ビットの配分としての不確定性関係}\label{ux534aux30d3ux30c3ux30c8ux306eux914dux5206ux3068ux3057ux3066ux306eux4e0dux78baux5b9aux6027ux95a2ux4fc2}

\textbf{著者：} 木原 範昭（ORCID: 0009-0004-6753-4020） \textbf{日付：}
2026年6月18日 \textbf{DOI（本版）：} 10.5281/zenodo.20742278 ／
\textbf{Concept DOI（全版）：} 10.5281/zenodo.20742277 ／
\textbf{Zenodo:} https://zenodo.org/records/20742278
\textbf{関連論文：}「高い対称性の極限におけるスケールの対数写像とスケール・量子化の分離について」§6.3（Concept
DOI: 10.5281/zenodo.20740841）

\begin{center}\rule{0.5\linewidth}{0.5pt}\end{center}

\subsection{趣旨}\label{ux8da3ux65e8}

本補遺は、関連論文
§6.3「揺らぎの保存則」で数行に圧縮した一点------\textbf{観測量空間の不確定性関係（積一定）が、基本量
N
空間における揺らぎの和ゼロから導出される}こと------を、独立に展開する。親論文に新しい主張を加えるものではなく、誤読されやすい導出の向きを明示することを目的とする。

特に強調するのは導出の\textbf{向き}である。通常、不確定性関係
δ\_a·δ\_b=一定（積一定）は基本法則として与えられる。本系ではこれは導出量であり、定義として置かれるのは
N 空間の\textbf{和ゼロ} Δ\_a+Δ\_b=0
の方である。和ゼロが定義、積一定が帰結------この向きが核心である。

本補遺は厳密な証明ではなく、親論文 §6.3
の構造を明示する観察である。揺らぎの大きさ ±1/2
は導出ではなく措定として扱う（親論文
§6.2）。物理量（位置・運動量等）への同定は行わない。ここでの「不確定性」は
Bell
の定理の意味ではなく、共役対に配分される量の保存という形式的構造に限る。

\begin{center}\rule{0.5\linewidth}{0.5pt}\end{center}

\subsection{1.
設定：基本量と揺らぎ}\label{ux8a2dux5b9aux57faux672cux91cfux3068ux63faux3089ux304e}

関連論文に従い、基本量を二乗量の対数として

\[N \equiv \tfrac{1}{2}\log_2(\nu^2) = \log_2\nu \qquad (N\in\mathbb{Z})\]

と置く。観測量 ν は逆写像 ν=2\^{}\{N\} で復元される。

共役な二量 a, b（観測量 ν, λ に対応）は、その積が保存される：

\[ab = k \quad\xrightarrow{\ \log_2\ }\quad N_a + N_b = K \qquad (K\equiv\log_2 k,\ \text{一定}).\]

すなわち N 空間では、共役対は\textbf{和の保存} N\_a+N\_b=K
を満たす。これは加法的恒等式であり、それ自体は構造を生まない（親論文 §5
の注）。構造は、この和の上に乗る揺らぎから生じる。

\begin{center}\rule{0.5\linewidth}{0.5pt}\end{center}

\subsection{2. 揺らぎの措定：半ビット
±1/2}\label{ux63faux3089ux304eux306eux63aaux5b9aux534aux30d3ux30c3ux30c8-12}

各基本量 N は揺らぎ Δ を持つ。本系ではこれを

\[\Delta = \pm\tfrac{1}{2} \qquad (\text{半ビット、措定、スケール不変})\]

と置く。この ±1/2 は導出ではなく措定である（親論文
§6.2：本系には標準偏差を構成する前提------観測の個別化・時間順序・値の持続------が無く、いかなる計算も隠れた前提を持ち込む）。揺らぎを観測量
ν の側でなく基本量 N の側に置くのは、N
空間で揺らぎを一定とすると観測量空間では揺らぎがスケール依存となり、整数計数の揺らぎとして整合的だからである。

「半ビット」と呼ぶのは、底 2 の対数空間における ±1/2 が、一ビット（N
の整数きざみ）の半分にあたるためである。底が 2 であることは親論文 §3・§4
による（系の二値性）。

\begin{center}\rule{0.5\linewidth}{0.5pt}\end{center}

\subsection{3.
定義：揺らぎの和ゼロ}\label{ux5b9aux7fa9ux63faux3089ux304eux306eux548cux30bcux30ed}

共役対の和は保存される（N\_a+N\_b=K、§1）。K
は定数であるから、両基本量の揺らぎ Δ\_a, Δ\_b
は、和を保つよう拘束される。これを本系の\textbf{定義}として置く：

\[\boxed{\ \Delta_a + \Delta_b = 0\ }\]

すなわち、一方が +1/2 揺れれば、他方は必ず −1/2
揺れる。半ビットを共役対のいずれに配分するかが揺らぎの自由度であり、その合計は常にゼロに保たれる。

この和ゼロは、共役対の和 N\_a+N\_b=K
が揺らぎを含めても保存される、という要請の言い換えである。揺らぎを加えた基本量を

\[\tilde N_a = N_a + \Delta_a, \qquad \tilde N_b = N_b + \Delta_b\]

と書けば、

\[\tilde N_a + \tilde N_b = (N_a+N_b) + (\Delta_a+\Delta_b) = K + 0 = K\]

となり、和は揺らぎの下でも厳密に K に保たれる。

幾何的には、共役対 (N\_a, N\_b) は直線 N\_a+N\_b=K 上の点であり、揺らぎ
Δ\_a=−1/2,
Δ\_b=+1/2（あるいはその逆）は、この直線に沿った滑りである。点は直線上に留まり、和は保存される（図1）。

\textbf{図1（fig1\_sum\_zero\_conserved\_line.svg）：} N
空間における和ゼロ。共役対 (N\_a, N\_b) は直線 N\_a+N\_b=K
上にあり、半ビットの揺らぎは一方を −1/2、他方を +1/2
ずらす。点は直線上を滑り、和は K に保たれる。

\begin{center}\rule{0.5\linewidth}{0.5pt}\end{center}

\subsection{4.
導出：積一定（不確定性関係）}\label{ux5c0eux51faux7a4dux4e00ux5b9aux4e0dux78baux5b9aux6027ux95a2ux4fc2}

和ゼロ（§3）を観測量空間に逆写像する。観測量は a=2\^{}\{\tilde N\_a\},
b=2\^{}\{\tilde N\_b\} で復元される：

\[ab = 2^{\tilde N_a + \tilde N_b} = 2^{K + 0} = 2^{K} \qquad (\text{一定}).\]

すなわち、和ゼロの拘束のもとで、観測量の積 ab は揺らぎに依らず一定
2\^{}K に保たれる。

観測量それぞれの揺らぎを考える。中心値 a\_0=2\^{}\{N\_a\} のまわりで、

\[a = 2^{N_a + \Delta_a} = a_0\cdot 2^{\Delta_a}, \qquad b = b_0\cdot 2^{\Delta_b}.\]

その積は

\[ab = a_0 b_0 \cdot 2^{\Delta_a + \Delta_b} = a_0 b_0 \cdot 2^{0} = a_0 b_0\]

であり、揺らぎ Δ\_a, Δ\_b がどう配分されても積は中心値の積 a\_0 b\_0
に固定される。和ゼロのもとでは、半ビットは各基本量に既約に ±1/2
残り、その符号が反相関する（一方が +1/2 なら他方は
−1/2）。観測量では一方が √2 倍なら他方は 1/√2 倍となり、積は a\_0 b\_0
に固定される。

これが観測量空間における\textbf{不確定性関係}の形である：半ビットは消せない床として各観測量に残り、共役対の積
a·b は a\_0
b\_0（=2\^{}K）に固定される。なお保たれるのは\textbf{観測量の積 a·b}
であって、\textbf{揺らぎの積 δ\_a·δ\_b ではない}（後者は配分 Δ
に依存し一定にならない）。Δ を連続に取る
squeeze（次段・図2）では、これは「一方を締めれば他方が広がる」トレードオフとして現れ、固定半ビット
±1/2 はその最小単位にあたる。この積一定は、N
空間の和ゼロ（§3、定義）から導出される帰結であって、独立に措定された法則ではない。

不確定性の最小単位は、措定された半ビット ±1/2
に対応する。なお単一の揺らぎ ±1/2
と、二つの揺らぎの積として現れる保存量（親論文 公理4 の
δ\textsubscript{1}δ\textsubscript{2}=1/2）は同一の量ではないが、同じ側------面積・シンプレクティック形式
ω（複素内積の虚部）------に属する量として対応づけられる。

観測量空間では、この保存関係は双曲線 a·b=2\^{}K
として現れる。一方を精密化（小さく）すれば他方が広がり、両者は双曲線上を移動するが積は一定に保たれる（図2）。N
空間の直線（和ゼロ）が、対数の逆写像により観測量空間の双曲線（積一定）に対応する。

\textbf{図2（fig2\_product\_constant\_observable.svg）：}
観測量空間における積一定。保存関係は双曲線 a·b=2\^{}K
であり、一方を締めれば他方が広がる。これは N
空間の和ゼロ（図1）の、対数逆写像による像である。

\begin{center}\rule{0.5\linewidth}{0.5pt}\end{center}

\subsection{5. 構造の要約}\label{ux69cbux9020ux306eux8981ux7d04}

{\def\LTcaptype{none} % do not increment counter
\begin{longtable}[]{@{}
  >{\raggedright\arraybackslash}p{(\linewidth - 6\tabcolsep) * \real{0.2500}}
  >{\raggedright\arraybackslash}p{(\linewidth - 6\tabcolsep) * \real{0.2500}}
  >{\raggedright\arraybackslash}p{(\linewidth - 6\tabcolsep) * \real{0.2500}}
  >{\raggedright\arraybackslash}p{(\linewidth - 6\tabcolsep) * \real{0.2500}}@{}}
\toprule\noalign{}
\begin{minipage}[b]{\linewidth}\raggedright
空間
\end{minipage} & \begin{minipage}[b]{\linewidth}\raggedright
量
\end{minipage} & \begin{minipage}[b]{\linewidth}\raggedright
揺らぎ
\end{minipage} & \begin{minipage}[b]{\linewidth}\raggedright
拘束
\end{minipage} \\
\midrule\noalign{}
\endhead
\bottomrule\noalign{}
\endlastfoot
基本量 N（定義側） & N=log\textsubscript{2}ν（整数） & Δ=±1/2（半ビット、措定） &
\textbf{Δ\_a+Δ\_b=0（和ゼロ、定義）} \\
観測量 ν（導出側） & ν=2\^{}N & δ=ν\textsubscript{0}(2\^{}Δ−1)（スケール依存） &
\textbf{a·b=2\^{}K（観測量の積が一定、導出）} \\
\end{longtable}
}

定義は N
空間の\textbf{和ゼロ}、導出は観測量空間の\textbf{積一定}である。不確定性の本体は「半ビット
±1/2
を共役対のいずれに配分するか、その合計は保存される」という配分の構造であり、積一定はその観測量空間での影である。

掛け算（積一定）と足し算（和ゼロ）は、対数写像 N=log\textsubscript{2}ν
で結ばれる。観測量空間で積として現れる不確定性が、基本量空間では和ゼロという単純な保存則になる------これが、対数を基本量に採ることの一つの帰結である。

\begin{center}\rule{0.5\linewidth}{0.5pt}\end{center}

\subsection{6.
限界（読み解きの注意）}\label{ux9650ux754cux8aadux307fux89e3ux304dux306eux6ce8ux610f}

本補遺は親論文 §6.3
の構造の明示であり、新しい主張を加えない。以下を明記する。

\begin{enumerate}
\def\labelenumi{\arabic{enumi}.}
\item
  \textbf{±1/2 は措定であり導出ではない。} 親論文 §6.2
  の通り、本系には標準偏差を構成する前提が無い。本補遺の導出（和ゼロ→積一定）は、この措定を前提とする。措定が外れれば帰結も変わる。
\item
  \textbf{和ゼロ Δ\_a+Δ\_b=0 は定義として置かれる。}
  これは共役対の和の保存（N\_a+N\_b=K）が揺らぎ下でも成り立つという要請の言い換えであり、より基本的な原理から導いたものではない。
\item
  \textbf{ここでの不確定性は Bell の定理の意味ではない。}
  共役対に配分される量の保存という形式的構造に限り、EPR/Bell の
  no-hidden-variables とは別物である。
\item
  \textbf{物理的同定を行わない。} 共役対 a, b（観測量 ν,
  λ）の物理的名称（位置・運動量・時間・エネルギー等）は無名のラベルであり、本補遺は特定の物理量への同定を行わない。
\item
  \textbf{§6.1 の整数きざみ揺らぎとは別物である。} 親論文 §6.1
  の「系の中の整数量が整数きざみ（±1）で揺らぐ」ことと、本補遺で N
  空間に置く半ビット揺らぎ Δ=±1/2
  は別個の措定であり、両者の関連（同一性・導出・近似）は主張しない。
\item
  \textbf{保存するのは観測量の積であり、揺らぎの積ではない。}
  和ゼロが保証するのは観測量の積 a·b=2\^{}K であって、揺らぎの積
  δ\_a·δ\_b
  ではない（後者は配分に依存し一定にならない）。また固定半ビットの和ゼロが与えるのは「各観測量に既約に残る
  ±1/2 の床＋符号の反相関」であり、Δ を連続に取る
  squeeze（双曲線上のトレードオフ「一方を締めれば他方が広がる」）とは別の側面である。±1/2
  はその squeeze の最小単位にあたる。両者を同一視しない。
\end{enumerate}

\begin{center}\rule{0.5\linewidth}{0.5pt}\end{center}

\subsection{参照}\label{ux53c2ux7167}

\begin{itemize}
\tightlist
\item
  Kihara,
  N.「高い対称性の極限におけるスケールの対数写像とスケール・量子化の分離について------最小公理系からの限定的検証報告（改訂版）」§6.2・§6.3。Concept
  DOI: 10.5281/zenodo.20740841。
\end{itemize}
\end{document}
